\chapter{Analyse des besoins et architecture du systeme}

\section{Introduction}
Ce chapitre presente les besoins metiers, les exigences techniques et l'architecture globale de la solution.

\section{Identification des acteurs}
\begin{itemize}
  \item \textbf{Touriste / Client} : consulte les informations hotelieres et interagit avec le chatbot.
  \item \textbf{Administrateur d'hotel} : met a jour les parametres, evenements et contenus.
  \item \textbf{Super administrateur} : supervise l'acces, la securite et les indicateurs globaux.
\end{itemize}

\section{Besoins fonctionnels}
\begin{itemize}
  \item authentification et verification des acces administrateurs ;
  \item consultation et edition des informations par hotel ;
  \item chatbot IA avec reponses contextuelles et multilingues ;
  \item streaming de reponses et conservation de contexte conversationnel ;
  \item collecte et visualisation des statistiques d'utilisation.
\end{itemize}

\section{Besoins non fonctionnels}
\begin{itemize}
  \item securite (JWT, validation des entrees, sanitation, rate limiting) ;
  \item performance (cache Redis, optimisation des routes API) ;
  \item maintenabilite (architecture modulaire) ;
  \item evolutivite (ajout de nouveaux hotels et modules) ;
  \item disponibilite et robustesse en production.
\end{itemize}

\section{Architecture generale}
L'architecture est basee sur Next.js (App Router) avec des API Routes cote serveur. Les donnees sont persistees dans PostgreSQL, le cache et les sessions sont geres via Redis, et la couche IA est assuree via un modele LLM integre dans le pipeline de chat.

\begin{figure}[H]
  \centering
  \begin{tikzpicture}[
    node distance=1.2cm and 1.4cm,
    box/.style={rectangle,draw,rounded corners,align=center,minimum width=3.2cm,minimum height=0.9cm},
    db/.style={cylinder,draw,shape border rotate=90,aspect=0.25,align=center,minimum height=1.1cm,minimum width=2cm},
    >={Latex}
  ]
    \node[box] (client) {Client\\(Interface Web)};
    \node[box,below=of client] (next) {Next.js 14\\Frontend + API Routes};
    \node[box,below left=of next,xshift=-0.5cm] (llm) {Service IA\\Groq LLM};
    \node[box,below=of next] (rag) {Base de connaissance\\RAG};
    \node[db,below right=of next,xshift=0.5cm] (pg) {PostgreSQL};
    \node[db,right=of pg,xshift=0.6cm] (redis) {Redis};
    \node[box,right=of next,xshift=1.0cm] (admin) {Dashboard Admin\\Analytics};

    \draw[->] (client) -- (next);
    \draw[->] (next) -- (llm);
    \draw[->] (next) -- (rag);
    \draw[->] (next) -- (pg);
    \draw[->] (next) -- (redis);
    \draw[->] (admin) -- (next);
  \end{tikzpicture}
  \caption{Architecture generale de la plateforme}
\end{figure}

\section{Technologies adoptees}
\begin{table}[H]
  \centering
  \caption{Technologies principales du projet}
  \begin{tabular}{p{4cm}p{4cm}p{6cm}}
    \toprule
    \textbf{Couche} & \textbf{Technologie} & \textbf{Role} \\
    \midrule
    Frontend & Next.js 14 + React 18 & Interface utilisateur et navigation \\
    Backend & API Routes Next.js & Logique serveur et endpoints REST \\
    Base de donnees & PostgreSQL & Stockage persistant \\
    Cache & Redis & Cache applicatif et sessions \\
    IA & Groq + LLM & Generation de reponses \\
    Visualisation & Recharts & Dashboards analytics \\
    \bottomrule
  \end{tabular}
\end{table}

\section{Conclusion}
Apres l'etude des besoins et de l'architecture cible, nous passons dans le chapitre suivant a la conception detaillee et a la modelisation des donnees.
